\documentclass{jsarticle}
\usepackage[dvipdfmx,final]{graphicx}
\usepackage{amsmath}
\usepackage{amssymb}
\usepackage{chemfig}
\begin{document}
\title{M theory in monopolity of \\
extra dimension\\
Symmetry theory}
\author{Masaaki Yamaguchi}
\date{}
\maketitle
   \newcommand{\variint}{%
	\begin{picture}(15,30)
		\put(0,0){
			$ \iint $}
		\put(0,5){\line(15,0){15}}
	\end{picture} %
}

\newcommand{\variiint}{%
	\begin{picture}(15,15)
		\put(0,0){
			$ \iiint $}
		\put(0,5){\line(15,0){20}}
	\end{picture} %
}



 \newcommand{\alearletter}{%
	 \begin{picture}(10,20)
		 \put(0,0){
			 $ \square $}
		 \put(0,0){\line(1,1){10}}
	 \end{picture} %
	 }

 \newcommand{\secproduct}{%
	 \begin{picture}(10,20)
		 \put(0,0){
			 $ \nabla $}
		 \put(5,0){+}
	 \end{picture} %
	 }
 \newcommand{\talot}{%
	 \begin{picture}(15,20)
		 \put(0,0){
			 $ \text{\scalebox{2.0}[2]{$\nabla$}}$}
		 \put(5,3){
			 $ \Delta $}
	 \end{picture} %
	 }
 \newcommand{\hazerdnabla}{%
	 \begin{picture}(15,20)
		 \put(0,0){
			 $ \text{\scalebox{2.0}[2]{$\nabla$}}$}
		 \put(5,3){
			 $\text{Y}  $}
	 \end{picture} %
	 }

 \newcommand{\securebox}{%
	 \begin{picture}(15,20)
		 \put(0,0){
			 $ \text{\scalebox{2.0}[2]{$\square$}}$}
		 \put(5,3){
			 $\text{Y}  $}
	 \end{picture} %
	 }


\newcommand{\timebow}{%
	\begin{picture}(5,5)
		\put(0,0){$\text{\rotatebox{90}{$\bowtie$}}$}
	\end{picture} %
	}


 \newcommand{\flowerletter}{%
	 \begin{picture}(10,20)
		 \put(0,0){$\text{\scalebox{1.0}[1]{$\square$}}$}
		 \put(0,0){$\text{\scalebox{1.0}[2]{$\int$}}$}
	 \end{picture} %
	 }
\newcommand{\squareiint}{%
	\begin{picture}(10,20)
		\put(0,0){
			$ \iint $}
		\put(0,4){$\text{\scalebox{3.0}[0.5]{$\square$}}$}
	\end{picture} %
}

\newcommand{\squareint}{%
	\begin{picture}(15,30)
		\put(0,0){
			$ \int $}
		\put(0,5){$\text{\scalebox{2.0}[1]{$\square$}}$}
	\end{picture} %
}


   \newcommand{\dazanier}{%
	\begin{picture}(12,12)
		\put(0,0){
			$ \bigsqcup $}
		\put(4,3){$\text{\scalebox{1.0}[1]{$\top$}}$}
		\put(4,3){$\text{\scalebox{1.0}[1]{$\bot$}}$}
	\end{picture} %
}
\newcommand{\innabla}{%
	\begin{picture}(13,13)
		\put(0,0){
			$ \nabla $}
		\put(0,0){$\text{\scalebox{2.0}[2]{$\nabla$}}$}
	\end{picture} %
}

\newcommand{\inndelta}{%
	\begin{picture}(20,20)
		\put(0,0){
			$ \Delta $}
		\put(0,0){$\text{\scalebox{2.0}[2]{$\Delta$}}$}
	\end{picture} %
}

\newcommand{\starnabla}{%
	\begin{picture}(20,20)
		\put(3,5){
			$ \Delta $}
		\put(0,0){$\text{\scalebox{2.0}[2]{$\nabla$}}$}
	\end{picture} %
}


  \newcommand{\whitewithbodercircle}
  {\text{\textcircled{\phantom{\textbullet}}}
  \kern-0.9em\text{\textcircled{\phantom{}}}}
     

 
  \newcommand{\whitewithblack}{\text{\textcircled{\textbullet}}}
  \newcommand{\whitewithblacksquare}{\text{\fbox{\text{\rule{0.5em}
  {0.5em}}}}}
  \newcommand{\whitewithwhitesquare}{\text{\fbox{\text{\phantom
  {\rule{0.5em}{0.5em}}}}}}
 \newcommand{\whitewithbodercircleiiint}{%
	 \begin{picture}(20,20)
		 \put(0,0){\scalebox{2.0}[2]{$\whitewithbodercircle$}}
		 \put(0,0){\scalebox{2.0}[2]{$\iiint$}}
	 \end{picture}%
	 }
\newcommand{\whitewithbodercircleint}{%
	 \begin{picture}(20,20)
		 \put(2,2){\scalebox{1.0}[1]{$\whitewithbodercircle$}}
		 \put(2,2){\scalebox{1.0}[2]{$\int$}}
	 \end{picture}%
	 }
\newcommand{\whitewithbodercircleiint}{%
	 \begin{picture}(20,20)
		 \put(0,0){\scalebox{1.0}[1]{$\whitewithbodercircle$}}
		 \put(2,2){\scalebox{1.0}[2]{$\iint$}}
	 \end{picture}%
	 }
 
 \newcommand{\whitewithbodercirclevarint}{%
	 \begin{picture}(13,13)
		 \put(0,0){\scalebox{2.0}[2]{$\whitewithbodercircle$}}
		 \put(0,0){\scalebox{2.0}[2]{$\varint$}}
	 \end{picture}%
	 }
 
 \newcommand{\whitewithbodercirclevariiint}{%
	 \begin{picture}(13,13)
		 \put(0,0){\scalebox{2.0}[2]{$\whitewithbodercircle$}}
		 \put(0,0){\scalebox{2.0}[2]{$\variiint$}}
	 \end{picture}%
	 }
 
 \newcommand{\average}{%
	  \begin{picture}(30,45)
		  \put(0,0){\scalebox{2.0}[2]{$\square$}}
		  \put(4,4){$\setminus$}
	  \end{picture}%
  } 
  \newcommand{\smallaverage}{%
	  \begin{picture}(30,45)
		  \put(0,0){\scalebox{1.0}[1]{$\square$}}
		  \put(0,0){$\setminus$}
	  \end{picture}%
  } 
 
  \newcommand{\whitewithboder}{%
	  \begin{picture}(15,30)
		  \put(0,0){\scalebox{2.0}[2]{$\square$}}
		  \put(4,4){$\square$}
	  \end{picture}%
  } 
  \newcommand{\whitewithnabla}{%
	  \begin{picture}(15,30)
		  \put(0,0){\scalebox{2.0}[2]{$\nabla$}}
		  \put(3,3){$\nabla$}
	  \end{picture}%
  } 


  The gut theory create with all of pond sentivelity
  from monopole of equation.

  統一場理論は、モノポールの磁気単極子である。
  $$ \int {\mathcal{K}(x,s)}^{\gamma((s,x),(x,0,1))} $$
  この方程式に導かれる式たちが以下である。
  $$ \pi^{e} \cong e^{\pi} $$
  $$ \pi^{\gamma(x,s)} \to \gamma \cong (\pi,e) $$
  サーストン・ペレルマン多様体が、以下である。
  $$ E(\sigma) = \mathcal{K}(\sigma) \otimes \mathcal{H}(\sigma) $$
  ガウス関数が、円周率となる結果である。
  $$ \int e^{-x^2 - y^2}dxdy = \pi $$
  $$ \int \pi^{\gamma(x,s)} = \int \int e^{\pi}dr $$
  重力場と反重力場の虚数のゼータ関数と実数のゼータ関数である。
  $$ {d \over {d f}}\int\int {1 \over {(x\log x)}^{2}}dx_m, $$
  $$ {d \over {d f}}\int\int {1 \over {(y \log y)}^{1 \over 2}}dy_m $$
  超越数がガンマ関数であることの、導かれる式が以下である。
  $$ \Gamma(s) = e = \log {\lim_{n\to \infty}{(1 + {1 \over n})}^{n}}  $$
  $$ \beta(p,q) = \pi = \int e^{\pi}dx $$
  以下、ガンマ関数とゼータ関数についての、ガンマ関数においての
  ベータ関数の結果である。
  $$ f(x) = e, f^{'}(x) = 0 $$
  $$ {f(x)}^{x} \cong x^{f(x)} $$
  $$ \Gamma^{'}(s) = 0, \gamma = \int e^{-x}x^{1 - t} \log x dx $$
  $$ = \int{{1 \over x^{s}} - \log x}dx = x $$
  $$ \int {\zeta(s)}^{\Gamma} = x $$
  Jones多項式が、ガンマ関数が、数を生み出す式である。
  $$ \int {\Gamma(s)}^{\Gamma} = x $$
  $$ \int ({{\zeta(s)} \over {x}})^{\Gamma} = 1 $$
  $$ \int {(\log x)}^{\Gamma} = 1 $$
  $$ \log x = {{\zeta(s)} \over {x}} $$
  シャノンの公式は、ゼータ関数である。
  $$ x\log x = \zeta(s) $$
  $$ \int {1 \over {x^s}}dx = \int x + \log x dx $$
  ガンマ関数は、超越数であり、
  $$ f(x) = e $$
  ゼータ関数は、オイラーの定数である。
  $$ f^{'}(x) = 0 $$
  ゼータ関数は、重力場と反重力場の積である。
  $$ \zeta(s) = {1 \over {\alearletter}} = {1 \over {square}} $$
  $$ \square = {1 \over {\alearletter}} = \square^{2} = 1 $$
  $$ \square \cdot \alearletter = 1 $$
  以下は、オイラーの定数が、ラムダ項になっている様相である。
  $$ \int\int e^{-x}x^{1 - t}\log x dx = 1 $$
  $$ \int\int e^{-x} x^{1 - t}dx_m = {1 \over {\log x}} $$
  $$ \int \Gamma dx = {1 \over {\log x}} $$
  $$ \log x \int \Gamma dx = 1 $$
  よって、ラムダ項がゼータ関数の結果である、宇宙の対称性となっている。
  $$ \sqrt{g} = 1 $$
   統一場理論は、モノポールのもう片方の磁気単極子でもある。
  $$ \int {\mathcal{H}(x,s)}^{\gamma((s,x),(x,0,1))} $$



\end{document}
